\section{ÍNDICE TENTATIVO}

Estructura de la Tabla de Contenidos proyectada para la Tesis Final, organizada según la normativa de la Carrera de Informática (UMSA) y el marco de \textit{Design Science Research} (DSR) de Vaishnavi y Kuechler (2004).

\begingroup
\renewcommand{\baselinestretch}{0.95}\selectfont
\footnotesize
\setlist[itemize]{noitemsep, topsep=0pt, parsep=0pt, partopsep=0pt}

\subsubsection*{CAPÍTULO I: MARCO INTRODUCTORIO}
\begin{itemize}
	\item[1.1.] Introducción
	\item[1.2.] Antecedentes
	\item[1.3.] Planteamiento del Problema
	\item[1.4.] Objetivos (General y Específicos)
	\item[1.5.] Hipótesis
	\item[1.6.] Justificación
	\item[1.7.] Alcances y Límites
	\item[1.8.] Aportes
\end{itemize}

\subsubsection*{CAPÍTULO II: MARCO TEÓRICO}
\begin{itemize}
	\item[2.1.] Votación Electrónica y Propiedades de Seguridad
	\item[2.2.] Criptografía de Clave Pública Clásica y su Vulnerabilidad Cuántica
	\item[2.3.] Criptografía Post-Cuántica Basada en Retículos (Module-LWE, Module-SIS, FIPS 203/204)
	\item[2.4.] Verificación Formal de Protocolos de Seguridad (Modelo Dolev--Yao, Tamarin Prover)
\end{itemize}

\subsubsection*{CAPÍTULO III: MARCO APLICATIVO}
\begin{itemize}
	\item[3.1.] Fase de Sugerencia: Formulación de la Propuesta del Protocolo
	\item[3.2.] Fase de Desarrollo: Fundamentación Teórica del Protocolo (Module-LWE / Module-SIS)
	\item[3.3.] Modelado Formal del Protocolo en Tamarin Prover
	\item[3.4.] Formalización de Lemmas de Confidencialidad, Integridad y Autenticidad
	\item[3.5.] Implementación de Referencia del Protocolo Propuesto
	\item[3.6.] Implementación del Protocolo Equivalente Basado en RSA/ECDSA
\end{itemize}

\subsubsection*{CAPÍTULO IV: PRUEBA DE HIPÓTESIS}
\begin{itemize}
	\item[4.1.] Fase de Evaluación: Verificación del Cumplimiento de los Lemmas Formales
	\item[4.2.] Medición de Indicadores de Rendimiento (Latencia, Tiempo de Ejecución, Consumo de Memoria)
	\item[4.3.] Comparación del Protocolo Propuesto frente al Protocolo Basado en RSA/ECDSA
	\item[4.4.] Contraste de la Hipótesis de Investigación
\end{itemize}

\subsubsection*{CAPÍTULO V: CONCLUSIONES Y RECOMENDACIONES}
\begin{itemize}
	\item[5.1.] Conclusiones
	\item[5.2.] Recomendaciones
\end{itemize}

\endgroup